% ==========================================================================
% Objetivos — se coloca en el front matter, justo después del Resumen.
% Estructura sin numeración de capítulo (chapter*/section*) pero visible
% en el índice mediante addcontentsline.
% ==========================================================================
\chapter*{Objetivos}
\label{cap:objetivos}
\phantomsection
\addcontentsline{toc}{chapter}{Objetivos}
\markboth{Objetivos}{}

\section*{Objetivo general}
\phantomsection
\addcontentsline{toc}{section}{Objetivo general}

Diseñar, implementar y validar un \emph{framework} bioinformático
reproducible que automatice la integración de cohortes transcriptómicas
públicas de cáncer de pulmón, mediante la combinación de un modelo de
lenguaje de gran tamaño para la curación clínica y aprendizaje automático
supervisado para la identificación de biomarcadores replicables.

El framework debe cumplir tres propiedades irrenunciables:

\begin{itemize}
  \item \textbf{Reproducibilidad}: dos ejecuciones sobre los mismos datos
  deben producir resultados idénticos hasta el último decimal, gracias al
  uso de semillas aleatorias fijas y a la ausencia de intervención manual
  en cualquier paso del pipeline.
  \item \textbf{Validación externa}: la evaluación del rendimiento debe
  realizarse siempre en cohortes independientes de las utilizadas para el
  entrenamiento, mediante Leave-One-Dataset-Out.
  \item \textbf{Interpretabilidad biológica}: los biomarcadores
  identificados deben ser trazables a mecanismos biológicos conocidos y,
  cuando exista un panel clínico establecido para la misma pregunta,
  contrastables con dicho panel.
\end{itemize}

\section*{Objetivos específicos}
\phantomsection
\addcontentsline{toc}{section}{Objetivos específicos}

Los objetivos específicos que se derivan del objetivo general son:

\begin{description}
  \item[OE1. Adquisición y normalización automatizada de cohortes GEO.]
  Descargar mediante API las cohortes seleccionadas del repositorio NCBI
  GEO, extraer las matrices de expresión y aplicar normalización log$_2$ y
  corrección por cuantiles dentro de cada estudio. El proceso debe ser
  robusto frente a la heterogeneidad de plataformas de microarray,
  particularmente Affymetrix HG-U133 Plus 2.0 (GPL570).

  \item[OE2. Curación clínica de metadatos con modelo de lenguaje.]
  Implementar un módulo que consulte al modelo Llama~3.3-70b a través de la
  API de Groq para asignar cada muestra a una categoría binaria
  (\texttt{sano} / \texttt{enfermo}), o descartarla como \texttt{ambigua},
  interpretando el texto libre de la descripción de la muestra. Registrar
  la tasa de éxito por cohorte como métrica de calidad de la curación.

  \item[OE3. Aplicación transparente de criterios de inclusión.] Definir y
  aplicar los criterios que una cohorte debe cumplir para entrar al
  entrenamiento y evaluación (presencia de casos y controles, curación
  exitosa por encima de un umbral mínimo, alineamiento verificado por
  identificador \texttt{geo\_accession}). Las cohortes que no cumplan
  algún criterio se documentan explícitamente, sin descartes silenciosos.

  \item[OE4. Análisis diferencial y clasificación supervisada.] Ejecutar
  análisis diferencial (test $t$ de Welch con corrección FDR de
  Benjamini-Hochberg) y entrenar clasificadores supervisados (regresión
  logística L1, Random Forest y máquinas de vectores soporte) para las dos
  tareas contempladas: (a) tumor \emph{vs.} sano, y (b) subtipo
  histológico ADC \emph{vs.} SQC.

  \item[OE5. Validación externa Leave-One-Dataset-Out.] Evaluar cada
  clasificador mediante LODO sobre las cohortes evaluables, reportando AUC,
  balanced accuracy, sensibilidad, especificidad y ganancia respecto al
  baseline mayoritario informado.

  \item[OE6. Identificación de una firma génica validada por consenso.]
  Construir una firma de biomarcadores mediante meta-análisis por consenso
  multi-cohorte: un gen entra en la firma sólo si mantiene el mismo signo
  de cambio y una magnitud mínima (medida como $d$ de Cohen) en las tres
  cohortes evaluables simultáneamente.

  \item[OE7. Panel mínimo clínicamente manejable.] Derivar, a partir de la
  firma validada, un panel mínimo del menor número de genes cuya AUC se
  mantenga a menos de 0,01 del máximo alcanzado por la firma completa. El
  objetivo es un panel del orden de las decenas de genes, susceptible de
  ser medido mediante PCR cuantitativa o \emph{NanoString}.

  \item[OE8. Validación externa contra marcadores clínicos IHC.]
  Comprobar en qué medida la firma validada, seleccionada de forma
  puramente estadística sin declararle los marcadores clínicos, recupera el
  panel de inmunohistoquímica diagnóstica recomendado por la OMS para
  adenocarcinoma y carcinoma escamoso.

  \item[OE9. Interfaz web reproducible.] Desarrollar una interfaz web
  interactiva (Streamlit) que permita a un revisor navegar por los
  resultados del framework, consultar cifras y figuras por cohorte y por
  gen, y verificar en tiempo real la trazabilidad de cualquier afirmación
  cuantitativa a los CSV/JSON generados por el pipeline.

  \item[OE10. Asistente conversacional acotado a los resultados.]
  Integrar en la interfaz un asistente conversacional basado en el mismo
  modelo Llama~3.3-70b, alimentado con un contexto que se lee dinámicamente
  de los resultados y con reglas explícitas que impiden inventar cifras o
  citar biomarcadores que no estén en la firma.
\end{description}

\section*{Hipótesis de partida}
\label{sec:hipotesis}
\phantomsection
\addcontentsline{toc}{section}{Hipótesis de partida}

Las hipótesis que se contrastan con el trabajo son:

\begin{enumerate}
  \item Un modelo de lenguaje contemporáneo puede curar los metadatos
  clínicos de cohortes GEO de cáncer de pulmón con una tasa de éxito
  superior al 80\,\% sobre cohortes evaluables, escalando la integración
  multi-cohorte a un coste computacional accesible.

  \item La aplicación de meta-análisis por consenso multi-cohorte, con
  exigencia estricta de replicación en tres cohortes independientes,
  produce una firma génica biológicamente coherente cuya AUC en LODO
  supera el 0,90 sobre la tarea de subtipo histológico.

  \item La firma resultante, seleccionada de forma \emph{data-driven},
  recupera al menos el 75\,\% de los marcadores diagnósticos IHC del panel
  recomendado por la OMS, lo que constituye una validación externa fuerte
  frente a un estándar clínico independiente.

  \item Un panel mínimo del orden de 20 genes es suficiente para
  aproximarse al rendimiento de la firma completa, con una pérdida de AUC
  inferior a 0,01, haciendo el panel clínicamente manejable.
\end{enumerate}

Los capítulos siguientes describen la metodología empleada
(capítulo~\ref{cap:metodos}) y presentan los resultados
(capítulo~\ref{cap:resultados}) que permiten aceptar o rechazar cada una
de estas hipótesis.
\cleardoublepage
